%% ==============================================================================
%% backmatter/glossary.tex
%% Glossaire des termes et table des acronymes du projet LaTeX
%% ==============================================================================
\addchap{Glossaire et Sigles}
\label{chap:glossaire}

\section*{Table des Sigles et Acronymes}
\addcontentsline{toc}{section}{Table des Sigles et Acronymes}

\noindent
\begin{tabularx}{\textwidth}{@{} >{\bfseries\color{ThesisPrimary}}p{2.6cm} X >{\raggedleft\arraybackslash\small\color{ThesisMuted}}p{3.5cm} @{}}
  \toprule
  Sigle & Définition / Signification complète & Domaine d'usage \\
  \midrule
  BCOR   & Binding Correction (Compensation de marge pour la reliure)    & Typographie KOMA \\
  BBX    & Bibliography Style File (Mise en page Bib\LaTeX{} des références) & Gestion bibliographique \\
  CBX    & Citation Style File (Mise en page des renvois dans le texte)  & Gestion bibliographique \\
  CLO    & Class Option File (Réglages d'options modulaires de classe)     & Moteur \LaTeX{} \\
  CLS    & Document Class File (Fichier racine de classe de document)     & Moteur \LaTeX{} \\
  CTAN   & Comprehensive \TeX{} Archive Network (Dépôt officiel des paquets) & Écosystème \TeX{} \\
  DIV    & Division Factor (Nombre de bandes pour la grille de page)      & Typographie KOMA \\
  DOI    & Digital Object Identifier (Identifiant pérenne d'article)      & Édition scientifique \\
  KOMA   & KOMA-Script (Ensemble de classes avancées par Markus Kohm)     & Typographie \LaTeX{} \\
  PDF    & Portable Document Format (Standard de diffusion vectoriel)     & Standard numérique \\
  STY    & Package Style File (Fichier d'extension ou de macros)           & Moteur \LaTeX{} \\
  TIKZ   & TikZ ist kein Zeichenprogramm (Moteur de dessin vectoriel)     & Graphisme \LaTeX{} \\
  TOC    & Table of Contents (Table des matières générale)                & Structure éditoriale \\
  \bottomrule
\end{tabularx}

\vspace{1.2cm}

\section*{Glossaire des Notions Typographiques}
\addcontentsline{toc}{section}{Glossaire des Notions Typographiques}

\setkomafont{descriptionlabel}{\bfseries\sffamily\color{ThesisPrimary}}
\begin{description}
  \item[Belle page] Dans un livre imprimé recto-verso, page impaire (située à droite lorsque l'ouvrage est ouvert),
  sur laquelle doivent traditionnellement s'ouvrir les têtes de chapitres (\code{open=right}).
  \item[Compensation de pliure (\code{BCOR})] Élargissement stratégique de la marge intérieure du bloc de texte,
  destiné à compenser l'espace absorbé physiquement par la reliure, l'encollage ou la couture de l'ouvrage.
  \item[Crénage] (\textit{kerning}) : ajustement millimétrique de l'espace optique séparant deux glyphes consécutifs
  (par exemple entre « V » et « A ») afin d'éviter les vides inesthétiques.
  \item[Découplage sémantique] Méthode d'ingénierie consistant à séparer les responsabilités
  d'un projet (orchestration, style, options d'affichage, contenu textuel pur).
  \item[Empagement (\textit{Satzspiegel})] Surface et proportions du rectangle de texte imprimé par rapport aux dimensions
  totales de la feuille de papier, calculé de façon optique par KOMA-Script via le diviseur \code{DIV}.
  \item[Microtypographie] Ensemble de techniques fines (protrusion des marges et expansion des polices)
  automatisées par le paquet \code{microtype} pour équilibrer la densité de gris typographique de la page.
  \item[Saut de ligne sémantique] Pratique de rédaction sous contrôle de version consistant
  à scinder le code source d'un paragraphe à la fin de chaque phrase ou proposition logique,
  rendant les historiques Git clairs et exempts de faux conflits.
\end{description}
